\section{The route toward quasi-polynomial recognition}
\label{sec:research}

\subsection{A complexity ledger}

The results in this article improve concrete operations. To determine what
they imply globally, one must retain both the encoded input size of each
operation and the number of operations. The following ledger records the
present status without replacing an unresolved factor by a constant.

\begin{table}[ht]
\centering\small
\begin{tabularx}{\linewidth}{p{0.23\linewidth}YY}
\toprule
Operation & Established here & Remaining global factor\\
\midrule
Full Jones computation & $\poly(n)2^{O(w)}$ bit time; certified
$w=O(\sqrt n)$ & Polynomial identity does not certify the unknot;
general polylogarithmic width is unavailable\\
Near-prefix compressed LCS & $O(gS^2)$ arithmetic operations;
polynomial bit cost for fixed $S$ & Search-move count, grammar growth,
and frequency of the small-deficit condition\\
Dihedral boundary transport & $O((k+1)B^3)$ bit time after preparation;
at most two root progressions & Number and encoding of cover
presentations; additional gluing data\\
Cyclic boundary marking types & At most
$\max\{2,-\chi(C)\}^{k-1}$ & Bounds on Euler complexity and attachment
slots for every generated piece\\
Hierarchy construction & Published compressed cutting is available &
Encoded-cost control of all construction, failure, repair, and iteration
branches\\
\bottomrule
\end{tabularx}
\caption{What each local result contributes to the target. A short
representation of a set does not by itself bound the number of sets
generated during a search.}
\label{tab:ledger}
\end{table}

\subsection{What an actual global proof must control}

Let $N(n)$ be the number of distinct algorithmic states actually processed,
let $s(n)$ bound the bit length of any state, and let $C(s)$ bound the bit
cost of processing a state, including certificate maintenance and canonical
comparison. An exhaustive but memoized implementation can have a bound
\begin{equation}\label{eq:global-ledger}
 T(n)\le\poly(n)+N(n)\,C(s(n))
\end{equation}
only after proving that all needed transitions are represented and that
the canonical identifications preserve the recognition problem. Hidden
enumeration of surfaces, gluing maps, or correction choices belongs in
$N$ or $C$; it cannot be omitted because individual objects have short
descriptions.

For the target $n^{O(\log n)}$, it would suffice to establish
$N(n)\le\exp(O(\log^2 n))$, $s(n)\le\exp(O(\log^2 n))$, and a fixed
polynomial bound on $C$. Stronger polynomial encoding bounds would make
such a proof more useful. These are sufficient accounting conditions,
not assertions that the current implementation satisfies them. In
particular, a polynomial certificate-verification theorem controls the
cost of checking a supplied object, but does not bound the search needed
to find it.

The iteration expression in Lackenby's Section~9 has the form
\begin{equation}\label{eq:lackenby-iterations}
 I\le L(G+1)^L,
\end{equation}
where $L$ bounds hierarchy length and $G$ bounds the relevant pattern
complexity. Its logarithm is at most
$\log L+L\log(G+1)$. Even if $G$ is polynomial in the original crossing
number, substituting only $L=O(n)$ into this estimate yields an
$\exp(O(n\log n))$ upper bound. It does not yield a quasi-polynomial
upper bound. This is a limitation of that inference, not a lower bound
on the algorithm. To obtain $n^{O(\log n)}$ from this form one needs,
for example, a proof that the \emph{effective} charged depth is
$O(\log n)$ with polynomial branching, or an amortized replacement that
controls the product of branch factors.

The distinction between existing compressed operations and unfinished
global construction is particularly important here. Lackenby's
Theorem~14.4 supplies polynomial compressed cutting for a compact properly
embedded orientable normal surface in an orientable 3-manifold with a
handle structure of uniform type. Proposition~14.2 additionally handles
admissibilization and removal of interior parallelity, assuming
irreducibility and an essential current pattern. Applying the latter on a branch whose
essentiality is still unknown requires a separate justification. The
research gap should be stated as encoded-cost and continuation control
for the entire algorithm, including failed candidates and repair, rather
than as the absence of a published compressed-cutting operation.

\subsection{A useful local consequence of the boundary theorem}

Suppose a fixed checked cyclic component $C$ has
$\max\{2,-\chi(C)\}\le n^a$ and $1\le k\le b\log(n+2)$ ordered
unparameterized boundary slots with their base boundary labels fixed,
where $a,b$ are constants. When $k=0$ there is one unmarked type. The
canonicalization theorem of Section~\ref{sec:transport} gives at most
\[
 (n^a)^{b\log(n+2)-1}=\exp(O(\log^2(n+2)))
\]
inequivalent markings of this component. The proved mixed-radix formula
also gives direct indexing of the canonical classes; the implemented
API returns a key, count, and witness, rather than running an enumerator.
Replacing a loop over all sheet
representatives by this canonical index is therefore a real local
quasi-polynomial reduction in the number of local marking classes,
provided the hypotheses hold. A corresponding running-time bound also
requires control of the presentation's bit size and its preparation
cost. Euler complexity alone does not control those costs: an annulus
can have an arbitrarily large binary sheet count while having only one
whole-circle marking class. The bound on class count is therefore kept
separate from the $B$-dependent arithmetic bounds.

There are three separate reasons this is not a general state-count
theorem. First, the component and covering presentation are fixed.
Second, point phases, parametrized attaching maps, orientation choices,
and interval-fibre prescriptions may contain additional information.
Third, different hierarchy levels can generate different pieces and
gluing graphs. One must bound the combination of these data, and prove
that discarded marking representatives have equivalent future behavior.

The exponential examples in Section~\ref{sec:transport} sharpen the
research task. Binary sheet count alone cannot replace Euler complexity.
Even genus zero does not control ordered point phases on an annulus.
Conversely, unparameterized boundary circles are substantially cheaper
than arbitrary point coordinates once the underlying topology is bounded.
A promising geometric implementation should therefore record the precise
kind of each attachment, rather than using point labels for every
boundary operation merely because points are convenient to serialize.

\subsection{Research questions for the tensor backend}

\paragraph{1. Minimize the turning cochain and integer size.}
The dual-tree construction proves a polynomial bound, but its $L_b$ need
not be small. For a fixed outer face, minimizing $\sum_e|b_e|$ subject to
the face equations is an integral minimum-cost flow problem on the dual
graph: replace each edge by two opposite arcs of unit cost and impose the
face demands. Integral demands admit an integral optimum; cancelling
opposite arc flow cannot increase cost. Thus a polynomial optimization
method can select a gauge with no larger $L_b$ than the tree solution.
Choosing the outer face as well can be reduced to trying all $n+2$ faces.
This observation supplies a concrete optimization problem; the current
patch uses the simpler tree construction and does not claim timings for
an unimplemented flow solver.

The practical question is whether smaller $L_b$ offsets the preparation
cost. A more refined objective would minimize the maximum accumulated
exponent span along the selected crossing order, rather than total
absolute turning. It must preserve the exact face equations and the
phase theorem. Compare tree, flow, and order-aware gauges on both the
successful grids and the regressing tree medials; record coefficient bit
lengths and total query time.

\paragraph{2. Choose between binary tensors and Potts partitions safely.}
The experiments show complementary strengths. Can a small bounded pilot
predict which representation will finish first without consuming a large
fraction of the total allowance? A useful portfolio would charge every
pilot transition, preserve common order certificates, and retain an
uncapped worst-case guarantee supplied by at least one branch. The
decision rule should be tested on held-out diagrams and with matched
orders. Merely choosing the fastest backend after seeing its completed
answer is not an implementable speedup.

An intermediate conversion is a separate question: can a processed Potts
summary be transformed exactly into the orientation tensor at a cost
smaller than recomputing the prefix? Such a map must account for turning
and the fixed specialization. A formula on arbitrary combinatorial
partitions cannot be inferred from a disk-boundary basis without checking
which partitions and phases are actually realizable.

\paragraph{3. Characterize useful families of small certified width.}
The binary theorem already gives $n^{O(\log n)}$ full Jones computation
on supplied orders of width $O(\log^2 n)$. Which knot-diagram families
encountered after the repository's structural simplifications admit such
orders, and can this be certified efficiently? General four-valent plane
graphs include large grids, so an argument based only on planarity cannot
force polylogarithmic width. The graph property must use additional
structure or a proved transformation of the diagram.

For nontrivial knots whose full polynomial is nonidentity, this can
accelerate the obstruction branch. A complete recognizer for a promised
family would additionally require a theorem that identity implies the
unknot in that family, or a separately bounded fallback. This promise
must be explicit in any claimed recognition complexity.

\paragraph{4. Find a homological analogue with an actual comparison map.}
The two-state tensor preserves an evaluated bracket, not a relative
chain complex. A research program toward Khovanov compression should
begin with a precisely defined retained object and a chain map, homotopy
equivalence, or sufficient detection theorem. What additional graded
data are needed when cancellations in the scalar tensor hide nonzero
homology? Does the required dimension grow polynomially, exponentially,
or faster with width? Small exhaustive counterexamples should be sought
before assuming a fixed-dimensional categorified boundary state.

This question is central rather than cosmetic: Khovanov detection gives
the fallback its mathematical completeness \cite{kronheimer-mrowka},
whereas the present Jones contraction only computes an obstruction.
A scalar identity cannot replace the missing homological theorem.

\subsection{Research questions for compressed group computations}

\paragraph{5. Prove a cyclic analogue of the near-prefix reduction.}
The new matcher is useful when a verified linear common prefix leaves
short tails. The current partial-relator operation may double words to
represent cyclic positions, making that deficit large. Can one work
directly with a long common cyclic arc and at most $S$ exceptional
letters, without doubling the relevant linear deficit? The desired
result should bound candidate cyclic alignments, preserve both donor
orientations, and produce a witness checked by the existing independent
relator verifier. Finding the long arc must be charged too.

\paragraph{6. Share periodic comparisons without losing exactness.}
Different near-prefix diagonals create related seeds and node--phase
tables. Can one share those tables across compatible periods or replace
them by a common refinement whose total size is controlled? Taking the
least common multiple of all periods can itself be expensive, so a
polynomial state bound must precede such a construction. Another
practical direction is an adaptive deficit gate based on reachable
grammar size and remaining work. The fixed $S\le32$ gate is deliberately
simple; its replacement should have an explicit charged bound and retain
the complete fallback.

\paragraph{7. Bound grammar growth and the number of certificate moves.}
Polynomial local word operations do not imply polynomial search. Prove
a potential that decreases by a controlled amount under \emph{batched}
Whitehead and relator moves while also bounding DAG growth. Expanded
length is unsuitable by itself: it can be exponentially large in the
input grammar, and unit decreases can be exponentially numerous.
The repository's existing power macros already remove particular
subtraction sequences. A general amortized theorem should track
symbol counts, grammar reuse, generator changes, and replay-certificate
size together, including branches in which no shortening move exists.

Evidence should include naturally arising relator states from difficult
diagrams where the new branch actually runs. The present fifteen-input
whole-knot audit records zero uses, so it supplies no empirical answer
to this question. Saving those future compressed states and their
independent replay witnesses would make the next experiment substantially
more informative than additional artificial powers alone.

\subsection{Research questions for geometric construction}

\paragraph{8. Turn compressed cutting into a complete checked interface.}
Implement the input and output objects of the published cutting theorem
with explicit boundary incidences and parallelity-bundle attachments.
Verify the output using compressed combinatorial identities or a separate
checker on bounded instances. Then account for every subsequent
admissibilization, compression, and repair operation when a candidate
fails an essentiality test. The key question is whether the same encoded
size bound survives all these operations, not whether a successful
essential hierarchy has a short certificate.

\paragraph{9. Prove a continuation theorem for canonical geometric states.}
Two pieces with the same genus and number of boundary components need
not admit the same gluings. Specify the complete retained data: fixed
base presentation, whole-circle slots, point or twist data when required,
orientations, interval fibres, and external incidence structure.
Prove that equality of canonical states gives mutually transportable
future operations and preserves the answer being decided. The cyclic
keys in this article solve one exact component of this specification.
They do not justify deleting the other data.

Extend the canonicalization from cyclic to the full supported dihedral
case, and then identify whether general AHT interval-pairing systems admit
a similarly small representation of all compatible maps. The at-most-two
progression theorem uses affine $\pm x+a$ actions in an essential way.
Piecewise affine systems may require more branches; their count and
bit size must be proved, not extrapolated from this special case.

\paragraph{10. Charge attachment complexity to topology across all levels.}
Can every piece produced by a suitable accelerated hierarchy be arranged
to have polynomial negative Euler complexity and only $O(\log n)$
simultaneously relevant unparameterized attachment slots? The local
marking theorem would then control that piece's marking choices.
The open issue is whether such bounds are preserved by cuts and repairs,
and whether the number of different underlying pieces is itself
quasi-polynomial. A counterexample with many indispensable attachment
slots would be just as valuable as a positive theorem, since it would
identify the need for a different summary.

\paragraph{11. Replace linear hierarchy depth by an amortized bound.}
A balanced divide-and-conquer decomposition can have logarithmic
recursion depth while still doing too much total work. Seek a potential
whose decreases control the \emph{product} of branching factors or the
number of distinct canonical states, including revisits and repairs.
One precise target is
$\sum_i\log b_i=O(\log^2 n)$ over the genuine branching steps of every
unfolded search path, together with polynomial bit work for each
intervening unary segment and a valid treatment of revisits. Unary work
must be charged separately: $2^n$ successive one-child updates have
branching product one but can still require exponential work.
Alternatively, a charged factor at least two for each nontrivial macro
update makes that cost visible in the logarithmic sum. A proof that only
the geometric height is logarithmic would leave the other factors in
\eqref{eq:global-ledger} unresolved.

\subsection{Research questions for certification and experiments}

\paragraph{12. Maintain a reproducible complexity and evidence ledger.}
For every new shortcut, record its theorem hypotheses, output guarantee,
encoded-size measure, resource semantics, and exact fallback. For every
performance claim, identify the completed operation, source revision,
inputs, warmups, paired controls, and censored runs. Expand the current
test corpus with independent exact answers and examples activating the
newly optimized path. Retain negative results and discovered ordering
defects; they indicate where the apparent bottleneck really lies.

A useful proof-assistant project is to formalize the finite identities
that guard the fastest paths: the face-cochain equations and their cycle
consequence, tensor expansion, balanced-base reconstruction, the
near-prefix coverage lemma, and the root-map/CRT solver. These are
well-scoped algebraic and combinatorial statements. Formalizing them
would strengthen implementation trust, while the global geometric
complexity theorem would remain a separate mathematical obligation.

\subsection{Suggested order of work}

The immediate implementation priorities are gauge-size reduction,
carefully charged backend selection, and recording real group-search
states that activate near-prefix matching. They have clear experimental
acceptance criteria and fit the current code.

The highest-value mathematical priorities for the quasi-polynomial
target are a complete compressed-cutting continuation interface,
canonical geometric state equivalence, and an amortized state-count
theorem. Work on the binary tensor's categorified analogue is a second
possible route, with the explicit warning that scalar cancellation
provides no automatic homological compression.

The present contribution reaches a stronger proved bound for full Jones
computation in this repository, removes an exact compressed-word
bottleneck, and supplies new explicit boundary-transport and local
counting tools. A general quasi-polynomial unknot-recognition theorem
would require the remaining global arguments in this section. None is
assumed in the implementation or inferred from the measurements.
